\documentclass[a4paper,12pt]{article}
\setcounter{secnumdepth}{5}
\setcounter{tocdepth}{3}
\input{/usr/share/LaTeX-ToolKit/template.tex}
\renewcommand{\arraystretch}{1.5}
\begin{document}
\title{計算機程式期末專題期末報告}
\author{沈威宇}
\date{\temtoday}
\titletocdoc
\sct{第一部分：專題整體}
\sct{第二部分：個人程式設計說明}

\ssc{報告架構}
在本章節中，我首先進行總論。此部分涵蓋要求的 A、B、C、D 四項題目，並進一步細分，闡述我在實作中所運用的 C++ 功能、概念、慣例以及最佳實踐。隨後，我將以模組或檔案群組為單位，依照課程要求，針對 A、B、C、D 四項題目（如適用）逐一進行說明。





\ssc{include/algorithm-copy.hpp}
The \verb|algorithm-copy.hpp| file provides three basic data transfer algorithms: \verb|copy|, \verb|copy_if|, and \verb|copy_backward|, mimicking the interface design of C++20 \verb|<algorithm>|. The copy function family demonstrates the polymorphism programming techniques, particularly conceptual parameter constraints and conditional optimization strategies.

I utilize concepts for constraints, such as \verb|std::input_iterator|, \verb|std::weakly_incrementable|, \verb|std::bidirectional_iterator|, and \verb|std::indirectly_writable| in the definition, and
\verb|std::contiguous_iterator| and \verb|std::is_trivially_copyable_v| for specialized acceleration.

In \verb|copy| algorithms, I use \verb|std::memcpy| when possible and element-wise assignment only when necessary, as demonstrated below (excerpted from \verb|include/algorithm-copy.hpp| lines 7-32)
\begin{lstlisting}[language=C++]
template <std::input_iterator InputIt, std::weakly_incrementable OutputIt>
  requires std::indirectly_writable<OutputIt, std::iter_value_t<InputIt>>
constexpr OutputIt copy(InputIt first, InputIt last, OutputIt d_first) {
  if constexpr (std::contiguous_iterator<InputIt> &&
                std::contiguous_iterator<OutputIt> &&
                std::is_trivially_copyable_v<
                    typename std::iterator_traits<InputIt>::value_type> &&
                std::is_same_v<
                    typename std::iterator_traits<InputIt>::value_type,
                    typename std::iterator_traits<OutputIt>::value_type>) {
    if (first == last)
      return d_first;
    const auto *src = std::to_address(first);
    const auto *src_end = std::to_address(last);
    auto *dest = std::to_address(d_first);
    const std::size_t count = (src_end - src) * sizeof(*src);
    std::memcpy(dest, src, count);
    return d_first + (last - first);
  }
  while (first != last) {
    *d_first = *first;
    ++first;
    ++d_first;
  }
  return d_first;
}
\end{lstlisting}

Similarly, in \verb|copy_backward| algorithm, \verb|std::memmove| is used when possible, as demonstrated below (excerpted from \verb|include/algorithm-copy.hpp| lines 49-73)
\begin{lstlisting}[language=C++]
template <std::bidirectional_iterator BidirIt1,
          std::bidirectional_iterator BidirIt2>
constexpr BidirIt2 copy_backward(BidirIt1 first, BidirIt1 last,
                                 BidirIt2 d_last) {
  using ValueType1 = typename std::iterator_traits<BidirIt1>::value_type;
  using ValueType2 = typename std::iterator_traits<BidirIt2>::value_type;
  if constexpr (std::contiguous_iterator<BidirIt1> &&
                std::contiguous_iterator<BidirIt2> &&
                std::is_trivially_copyable_v<ValueType1> &&
                std::is_same_v<ValueType1, ValueType2>) {
    auto n = last - first;
    if (n > 0) {
      std::memmove(std::addressof(*d_last) - n, std::addressof(*first),
                   n * sizeof(ValueType1));
    }
    return d_last - n;
  } else {
    while (first != last) {
      --last;
      --d_last;
      *d_last = *last;
    }
    return d_last;
  }
}
\end{lstlisting}








I originally struggled on implementating \verb|operator<=>| for \verb|_tuple_impl|, the internal recursive inheritance implementation for tuple, since different element may return different type for auto deduct to fail and deducing type recursively with conditional_t will face initialization of undefined template parameters even though it will not be chose. Even though it succeeded to compile in g++, that is not complying with standard and not compiled in clang. Latter, I add element_type internal type alias in \verb|_tuple_impl| that calculate based on index and recursively define the alias. And use it and get in tuple directly to deduce type at the initialization iteratively instead of recursively, which succeeded.
